% Theory --- formalization spine.
% Every theorem and corollary here has a corresponding entry in
% notes/formalization_status.csv and notes/theorem_cards/T*.md.
% Proofs are written to match the Lean targets listed alongside each
% statement (via \leanref{...}), even when the Lean file is still in
% progress; the `status` field in the CSV tracks which are machine-checked.

In this section we collect the theoretical results that motivate
and justify the LEMO architecture.  The development is in three
parts.  First, on the discrete cyclic group $C_n$, we establish the
equivariance spine of the architecture and derive the finite-group
out-of-distribution (OOD) gap that motivates shift-augmentation.
Second, we introduce the spectrally-normalized subclass $\mathrm{LEMO}_\sigma$
and prove an architectural contraction bound that certifies rollout
stability.  Third, we extend the equivariance story to the continuous
circle $S^1$, where no finite augmentation can match the equivariant
learner.

\subsection{Setup: the cyclic shift action and the lag convolution}

Throughout this section, fix an integer $n \ge 1$ and work over the
cyclic group $C_n = \mathbb{Z}/n\mathbb{Z}$.

\begin{definition}[History space and shift action]
\label{def:shift}
The \emph{history space} is $\mathcal{X}_n := \{h : C_n \to \mathbb{R}\}$.
For each $r \in C_n$, the \emph{cyclic shift} $\mathrm{sh}_r : \mathcal{X}_n \to \mathcal{X}_n$
is defined by $(\mathrm{sh}_r h)(j) := h(j - r)$.
\end{definition}

\begin{theorem}[Shift action]
\label{thm:shift-action}
The family $\{\mathrm{sh}_r\}_{r \in C_n}$ is a group action of $C_n$ on
$\mathcal{X}_n$: $\mathrm{sh}_0 = \mathrm{id}$ and $\mathrm{sh}_{r+s} = \mathrm{sh}_r \circ \mathrm{sh}_s$ for all
$r, s \in C_n$.
\end{theorem}

\begin{proof}
Direct: $(\mathrm{sh}_0 h)(j) = h(j - 0) = h(j)$, and
$(\mathrm{sh}_{r+s} h)(j) = h(j - r - s) = (\mathrm{sh}_r (\mathrm{sh}_s h))(j)$ by
associativity of subtraction in $C_n$.
\leanref{LEMO/DiscreteCore.lean}{shift\_is\_action}
\end{proof}

\begin{definition}[Lag convolution]
\label{def:lagconv}
For a kernel $K : C_n \to \mathbb{R}$, the \emph{cyclic lag-convolution
operator} $\mathrm{L}_K : \mathcal{X}_n \to \mathcal{X}_n$ is
\[
  (\mathrm{L}_K h)(j) := \sum_{k \in C_n} K(k) \cdot h(j - k).
\]
\end{definition}

\begin{theorem}[Lag-convolution equivariance]
\label{thm:layer-equiv}
For every kernel $K : C_n \to \mathbb{R}$ and every $r \in C_n$,
\[
  \mathrm{L}_K \circ \mathrm{sh}_r \;=\; \mathrm{sh}_r \circ \mathrm{L}_K.
\]
\end{theorem}

\begin{proof}
For any $h \in \mathcal{X}_n$ and $j \in C_n$,
\[
  (\mathrm{L}_K (\mathrm{sh}_r h))(j)
  = \sum_{k} K(k) \cdot h(j - r - k)
  = \sum_{k} K(k) \cdot h((j - r) - k)
  = (\mathrm{L}_K h)(j - r)
  = (\mathrm{sh}_r (\mathrm{L}_K h))(j).
\]
The middle equality is associativity of subtraction in $C_n$.
\leanref{LEMO/DiscreteCore.lean}{lagConv\_shift}
\end{proof}

\begin{theorem}[Stack equivariance]
\label{thm:stack-equiv}
Let $f_1, \dots, f_L : \mathcal{X}_n \to \mathcal{X}_n$ be shift-equivariant
maps, i.e., $f_i \circ \mathrm{sh}_r = \mathrm{sh}_r \circ f_i$ for every $i$ and every
$r \in C_n$.  Then their composition $f_L \circ \cdots \circ f_1$ is
shift-equivariant.
\end{theorem}

\begin{proof}
Induction on $L$.  The base case $L = 1$ is the hypothesis.  For the
step, $(f_{L+1} \circ f_L \circ \cdots \circ f_1) \circ \mathrm{sh}_r =
f_{L+1} \circ (f_L \circ \cdots \circ f_1 \circ \mathrm{sh}_r) =
f_{L+1} \circ \mathrm{sh}_r \circ (f_L \circ \cdots \circ f_1) =
\mathrm{sh}_r \circ f_{L+1} \circ \cdots \circ f_1$.
\leanref{LEMO/DiscreteCore.lean}{stack\_equivariant}
\end{proof}

\begin{theorem}[Softmax equivariance]
\label{thm:softmax-equiv}
For every temperature $\tau > 0$, the normalized softmax along the lag
axis,
\[
  (\mathrm{softmax}_\tau h)(j) := \frac{\exp(h(j)/\tau)}{\sum_{k \in C_n} \exp(h(k)/\tau)},
\]
commutes with every cyclic shift: $\mathrm{softmax}_\tau \circ \mathrm{sh}_r = \mathrm{sh}_r \circ \mathrm{softmax}_\tau$.
\end{theorem}

\begin{proof}
The denominator $Z(h) := \sum_k \exp(h(k)/\tau)$ is shift-invariant:
$Z(\mathrm{sh}_r h) = \sum_k \exp(h(k-r)/\tau) = \sum_{k'} \exp(h(k')/\tau) = Z(h)$
by reindexing $k' = k - r$ (which is a bijection on $C_n$).  Hence
$(\mathrm{softmax}_\tau (\mathrm{sh}_r h))(j) = \exp(h(j-r)/\tau)/Z(h) =
(\mathrm{softmax}_\tau h)(j - r) = (\mathrm{sh}_r (\mathrm{softmax}_\tau h))(j)$.
\leanref{LEMO/DiscreteCore.lean}{softmax\_shift}
\end{proof}

\begin{theorem}[Pooled invariance]
\label{thm:pooled-inv}
The sum-pool $\mathrm{sumpool}(h) := \sum_{j \in C_n} h(j)$ is
shift-invariant: $\mathrm{sumpool}(\mathrm{sh}_r h) = \mathrm{sumpool}(h)$
for every $r \in C_n$.  The same holds for mean-pool,
$\mathrm{logsumexp}$, and any symmetric function of the lag axis.
\end{theorem}

\begin{proof}
$\mathrm{sumpool}(\mathrm{sh}_r h) = \sum_j h(j-r) = \sum_{j'} h(j') = \mathrm{sumpool}(h)$
by the reindexing $j' = j - r$.
\leanref{LEMO/Pooling.lean}{sumpool\_invariant}
\end{proof}

\begin{lemma}[Orbit constancy]
\label{lem:orbit-constancy}
Let $f : \mathcal{X}_n \to \mathcal{X}_n$ be shift-equivariant and let $\mathcal{L} :
\mathcal{X}_n \times \mathcal{X}_n \to \mathbb{R}_{\ge 0}$ be shift-invariant (i.e.,
$\mathcal{L}(\mathrm{sh}_r y, \mathrm{sh}_r y') = \mathcal{L}(y, y')$ for every $r$).
Then for every $h, h' \in \mathcal{X}_n$ and every $r \in C_n$,
\[
  \mathcal{L}\bigl(f(\mathrm{sh}_r h), f(\mathrm{sh}_r h')\bigr)
  \;=\; \mathcal{L}\bigl(f(h), f(h')\bigr).
\]
In particular, for any fixed equivariant target $f^\star$, the
function $h \mapsto \mathcal{L}(f(h), f^\star(h))$ is constant on
$C_n$-orbits.
\end{lemma}

\begin{proof}
Substitute equivariance of $f$ and invariance of $\mathcal{L}$:
$\mathcal{L}(f(\mathrm{sh}_r h), f(\mathrm{sh}_r h')) =
\mathcal{L}(\mathrm{sh}_r f(h), \mathrm{sh}_r f(h')) = \mathcal{L}(f(h), f(h'))$.
\leanref{LEMO/Pooling.lean}{orbit\_constant}
\end{proof}

\subsection{Finite-group OOD gap}

Let $X$ be a finite set on which $C_n$ acts freely, so that $|X| = n k$
where $k$ is the number of orbits.  Fix a \emph{fundamental domain}
$D_0 \subset X$ containing exactly one element of each orbit; let
$\mu_\mathrm{fund}$ be the uniform measure on $D_0$ and $\mu_\mathrm{uni}$
the uniform measure on $X$ (equivalently, the orbit-average of
$\mu_\mathrm{fund}$).  Let $Y = \{0, 1\}$ and consider the $0/1$ loss
$\ell(y, y') := \mathbf{1}[y \ne y']$.

\begin{theorem}[Finite-group OOD gap]
\label{thm:oodgap}
Let $f^\star : X \to Y$ be shift-invariant.
\begin{enumerate}
  \item[(i)] For every shift-invariant hypothesis $\hat{f} : X \to Y$,
  $\mathbb{E}_{\mu_\mathrm{fund}}[\ell(\hat{f}, f^\star)] =
  \mathbb{E}_{\mu_\mathrm{uni}}[\ell(\hat{f}, f^\star)]$.
  \item[(ii)] There exists a (non-invariant) hypothesis
  $\hat{f} : X \to Y$ with zero fundamental-domain risk but full-orbit
  risk exactly $(n-1)/n$.
\end{enumerate}
\end{theorem}

\begin{proof}
\textbf{(i)}~By Lemma~\ref{lem:orbit-constancy}, $x \mapsto \ell(\hat{f}(x),
f^\star(x))$ is constant on $C_n$-orbits.  The orbit-average of a
function that is constant on orbits equals its value on any orbit
representative; averaging over $D_0$ gives the same value as averaging
over the full orbit-average $\mu_\mathrm{uni}$.

\textbf{(ii)}~Define $\hat{f}(x) := f^\star(x)$ if $x \in D_0$ and
$\hat{f}(x) := 1 - f^\star(x)$ otherwise.  Then $\hat{f}$ has zero
fundamental-domain risk; on any orbit, $\hat{f}$ disagrees with $f^\star$
at $n - 1$ of the $n$ points, so the full-orbit risk is exactly
$(n-1)/n$.
\leanref{LEMO/OODGap.lean}{finite\_group\_ood\_gap}
\end{proof}

\begin{corollary}[PAC sample-complexity separation]
\label{cor:pac-gap}
In the setting of Theorem~\ref{thm:oodgap}, let $\mathcal{F}_G$ be the
class of shift-invariant functions $X \to Y$ (with
$|\mathcal{F}_G| = 2^k$) and $\mathcal{F} = Y^X$ the full class (with
$|\mathcal{F}| = 2^{nk}$).  To learn a target in $\mathcal{F}_G$ under
$\mu_\mathrm{uni}$ to excess $0/1$ risk $\varepsilon$ with probability
$1 - \delta$, empirical risk minimization needs
\[
  m_\mathrm{eq} \;=\; \Theta\!\left(\frac{k \log(1/\varepsilon) + \log(1/\delta)}{\varepsilon}\right),
  \qquad
  m_\mathrm{non} \;=\; \Theta\!\left(\frac{nk \log(1/\varepsilon) + \log(1/\delta)}{\varepsilon}\right),
\]
respectively, so the ratio $m_\mathrm{non}/m_\mathrm{eq} = \Theta(n)$ for
the leading term.
\end{corollary}

\begin{proof}
Upper bounds are the standard finite-hypothesis-class PAC bounds for
$\mathcal{F}_G$ and $\mathcal{F}$.  The matching lower bound follows
from the construction of Theorem~\ref{thm:oodgap}(ii): there is a
subfamily of $\mathcal{F}$ of size $2^{(n-1)k}$ that is perfectly
distinguishable on $\mu_\mathrm{uni}$ but indistinguishable on the
fundamental-domain marginal, so sample-efficient learning on
$\mu_\mathrm{uni}$ from fundamental-domain-only data is impossible
without the invariance prior.  The construction follows
Mei--Misiakiewicz--Montanari (2021) and Elesedy--Zaidi (2021).
\leanref{LEMO/OODGap.lean}{pac\_sample\_complexity}
\end{proof}

\subsection{Operator-norm control}

We now turn to the ingredients for rollout stability.  All three
lemmas below are standard; we collect them because they combine into
the architectural contraction bound of
Corollary~\ref{cor:architectural-contraction}.

\begin{lemma}[FFT norm bound]
\label{lem:fft-norm}
For $K : C_n \to \mathbb{R}$, let $\hat{K} : C_n \to \mathbb{C}$,
$\hat{K}(\omega) := \sum_{k \in C_n} K(k) \exp(-2\pi i k \omega / n)$ be
its discrete Fourier transform.  Viewing $\mathrm{L}_K$ as an operator on
$\ell^2(C_n)$, its Euclidean operator norm is
\[
  \|\mathrm{L}_K\|_\mathrm{op}
  \;=\; \max_{\omega \in C_n} |\hat{K}(\omega)|.
\]
\end{lemma}

\begin{proof}
The DFT is unitary on $\ell^2(C_n)$.  Under the DFT basis, $\mathrm{L}_K$ is
the diagonal operator with entries $\hat{K}(\omega)$.  The operator
norm of a diagonal operator is the maximum magnitude of its diagonal.
\leanref{LEMO/FFTNorm.lean}{fft\_op\_norm}
\end{proof}

\begin{lemma}[Pointwise 1-Lipschitz operator bound]
\label{lem:pointwise-norm}
Let $\sigma : \mathbb{R} \to \mathbb{R}$ be $1$-Lipschitz.  Extend
$\sigma$ pointwise to $\mathcal{X}_n^d := (\mathbb{R}^d)^n$.  Then for every
$H, H' \in \mathcal{X}_n^d$,
\[
  \|\sigma(H) - \sigma(H')\|_2 \;\le\; \|H - H'\|_2.
\]
\end{lemma}

\begin{proof}
$\|\sigma(H) - \sigma(H')\|_2^2 = \sum_{j, c} |\sigma(H_{j,c}) - \sigma(H'_{j,c})|^2
\le \sum_{j, c} |H_{j,c} - H'_{j,c}|^2 = \|H - H'\|_2^2$ by
$1$-Lipschitzness at each coordinate.
\leanref{LEMO/FFTNorm.lean}{pointwise\_op\_norm}
\end{proof}

\begin{lemma}[Composition Lipschitz]
\label{lem:composition-lip}
If $f : M \to N$ is $L_1$-Lipschitz and $g : N \to P$ is $L_2$-Lipschitz
between metric spaces, then $g \circ f : M \to P$ is
$(L_1 L_2)$-Lipschitz.
\end{lemma}

\begin{proof}
$d_P(g(f(x)), g(f(y))) \le L_2 \cdot d_N(f(x), f(y)) \le L_2 L_1 \cdot d_M(x, y)$.
\leanref{LEMO/FFTNorm.lean}{composition\_lipschitz}
\end{proof}

\subsection{The $\mathrm{LEMO}_\sigma$ subclass: architectural contraction}

\begin{definition}[$\mathrm{LEMO}$ predictor and the $\mathrm{LEMO}_\sigma$ subclass]
\label{def:lemo-sigma}
A \emph{$\mathrm{LEMO}$ predictor} is a map $\Phi : (\mathbb{R}^d)^n \to \mathbb{R}^d$
of the form $\Phi = \pi \circ \sigma_L \circ W_L \circ \cdots \circ
\sigma_1 \circ W_1$, where each $W_\ell$ is a lag-convolution layer
(Definition~\ref{def:lagconv}), each $\sigma_\ell$ is a pointwise
$1$-Lipschitz nonlinearity, and $\pi$ is a shift-equivariant readout.
For $\sigma \in (0, 1)$, the subclass $\mathrm{LEMO}_\sigma$ consists of
all $\mathrm{LEMO}$ predictors in which every $W_\ell$ and $\pi$ has Euclidean
operator norm at most $\sigma$.
\end{definition}

In practice, the per-layer constraint $\|W_\ell\|_\mathrm{op} \le \sigma$
for convolutional layers is enforced by dividing each Fourier coefficient
$\hat{W_\ell}(\omega)$ by $\max(1, \max_\omega |\hat{W_\ell}(\omega)|/\sigma)$
at every forward pass (cf.~Lemma~\ref{lem:fft-norm}).  Importantly,
this is \emph{not} what
\texttt{torch.nn.utils.spectral\_norm(nn.Conv1d)} computes; that
wrapper bounds a reshaped-weight SVD and is typically looser (and sometimes
arbitrarily loose) than the convolutional operator norm.  The FFT-based
normalization is implemented in \texttt{src/models/lemo.py}.

\begin{corollary}[Architectural contraction]
\label{cor:architectural-contraction}
Every $\Phi \in \mathrm{LEMO}_\sigma$ is $\sigma^{L+1}$-Lipschitz in the
Euclidean norm:
\[
  \|\Phi(H) - \Phi(H')\|_2 \;\le\; \sigma^{L+1} \cdot \|H - H'\|_2
  \qquad \forall\, H, H' \in (\mathbb{R}^d)^n.
\]
\end{corollary}

\begin{proof}
By Lemma~\ref{lem:fft-norm} each $W_\ell$ is $\sigma$-Lipschitz in the
Euclidean norm, and the readout $\pi$ is $\sigma$-Lipschitz by the same
argument (or by direct construction).  Each $\sigma_\ell$ is
$1$-Lipschitz by Lemma~\ref{lem:pointwise-norm}.  The composition is
$(\sigma \cdot 1)^L \cdot \sigma = \sigma^{L+1}$-Lipschitz by
Lemma~\ref{lem:composition-lip} applied $2L + 1$ times.
\leanref{LEMO/Contraction.lean}{lemo\_sigma\_contract}
\end{proof}

\begin{corollary}[Rollout perturbation]
\label{cor:rollout}
Let $\Phi \in \mathrm{LEMO}_\sigma$ with depth $L$ and history length $n$,
and let $U_\Phi$ be the associated history-update map.  There exists
$\beta > 1$ (specifically $\beta = \sigma^{-(L+1)/n}$) such that $U_\Phi$
is $\rho$-contractive in the weighted norm
$|H|_\beta := \max_j \beta^{-j} |H_j|$ with rate
$\rho = \sigma^{(L+1)/n} < 1$.  Consequently, if the one-step prediction
error satisfies $|\hat{\Phi}(H_t) - \Phi(H_t)|_\beta \le \varepsilon$ for
all $t$, then the rollout error is bounded by
\[
  |H_t - \hat{H}_t|_\beta
  \;\le\; \frac{1 - \rho^t}{1 - \rho} \cdot \varepsilon
  \;\le\; \frac{\varepsilon}{1 - \rho}
  \qquad \forall\, t \ge 0.
\]
\end{corollary}

\begin{proof}
By Corollary~\ref{cor:architectural-contraction}, $\Phi$ is
$\sigma^{L+1}$-Lipschitz in the Euclidean norm.  A standard
comparability estimate between the Euclidean and weighted norms on
$(\mathbb{R}^d)^n$ gives $\|\cdot\|_2$--$|\cdot|_\beta$ bi-Lipschitz
constants that absorb the $\beta^{n-1}$ factor when $\beta$ is chosen
as in the statement; substituting into the Euclidean Lipschitz bound
yields $\rho = \sigma^{(L+1)/n}$.  The rollout bound then follows
by induction: if $E_t := |H_t - \hat{H}_t|_\beta$ satisfies
$E_{t+1} \le \rho E_t + \varepsilon$, then
$E_t \le (1 - \rho^t)/(1 - \rho) \cdot \varepsilon$.
\leanref{LEMO/Contraction.lean}{rollout\_bound}
\end{proof}

\subsection{Continuous-lag extension}
\label{sec:continuous-lag}

The discrete theory above uses shifts in $C_n$.  In many applications
(including DDEs) the lag is a continuous quantity $\tau \in [0, L]$.
Finite augmentation over discrete shifts cannot fill the continuum,
so the OOD story sharpens considerably.  Throughout this subsection
let $\mathcal{X} := L^\infty(S^1; \mathbb{R})$ with $S^1 = \mathbb{R}/L\mathbb{Z}$ and let
$\mathrm{sh}_s : \mathcal{X} \to \mathcal{X}$, $(\mathrm{sh}_s h)(\tau) := h((\tau - s) \bmod L)$, be
the continuous shift action.

\begin{definition}[Continuous lag-convolution operator]
\label{def:continuous-lagconv}
For $K \in L^1(S^1)$ the continuous lag-convolution
$\mathrm{L}_K : \mathcal{X} \to \mathcal{X}$ is
\[
  (\mathrm{L}_K h)(\tau) := \int_0^L K\bigl((\tau - s) \bmod L\bigr) \cdot h(s)\,ds.
\]
\end{definition}

\begin{theorem}[Continuous-lag equivariance]
\label{thm:continuous-layer-equiv}
For every $K \in L^1(S^1)$ and every $s \in S^1$,
$\mathrm{L}_K \circ \mathrm{sh}_s = \mathrm{sh}_s \circ \mathrm{L}_K$.
\end{theorem}

\begin{proof}
Translation-invariance of Lebesgue measure on $S^1$:
\[
  (\mathrm{L}_K (\mathrm{sh}_s h))(\tau)
  = \int K((\tau - u) \bmod L) \cdot h((u - s) \bmod L)\, du
  = \int K(((\tau - s) - v) \bmod L) \cdot h(v)\, dv
  = (\mathrm{sh}_s (\mathrm{L}_K h))(\tau),
\]
substituting $v = (u - s) \bmod L$.
\leanref{LEMO/ContinuousQuadrature.lean}{continuous\_lag\_equiv\_discrete}
\end{proof}

\begin{theorem}[Continuous OOD gap]
\label{thm:continuous-ood-gap}
Fix a fundamental-domain measure $\mu_\mathrm{fund} \in \mathcal{P}(D_0)$ for
some $D_0 \subset \mathcal{X}$ meeting every $S^1$-orbit exactly once, and
let $\mu_\mathrm{uni}$ be its orbit-average (the pushforward of
$\mu_\mathrm{fund} \otimes \mathrm{Unif}(S^1)$).  Let $f^\star : \mathcal{X} \to \mathcal{X}$
be continuous-lag-equivariant.
\begin{enumerate}
\item[(i)] For every continuous-lag-equivariant hypothesis $\hat{f}$,
$\mathbb{E}_{\mu_\mathrm{fund}}[\mathcal{L}(\hat{f}, f^\star)] = \mathbb{E}_{\mu_\mathrm{uni}}[\mathcal{L}(\hat{f}, f^\star)]$.
\item[(ii)] For binary output and $0/1$ loss, the non-equivariant
construction of Theorem~\ref{thm:oodgap}(ii) yields a hypothesis
$\hat{f}$ with zero fundamental-domain risk and full-orbit risk
$1 - |D_0|/|\mathcal{X}|_\mathrm{uni}$, which tends to $1$ as the group becomes
continuous.
\end{enumerate}
\end{theorem}

\begin{proof}
\textbf{(i)} By Theorem~\ref{thm:continuous-layer-equiv} the loss
$h \mapsto \mathcal{L}(\hat{f}(h), f^\star(h))$ is constant on orbits; integrals of
orbit-constant functions agree under $\mu_\mathrm{fund}$ and its
orbit-average $\mu_\mathrm{uni}$.
\textbf{(ii)} Same construction as Theorem~\ref{thm:oodgap}(ii), now
applied to the continuous orbit structure.
\leanref{LEMO/ContinuousQuadrature.lean}{continuous\_ood\_gap\_discrete}
\end{proof}

\begin{theorem}[Augmentation covering-radius lower bound]
\label{thm:aug-lower-bound}
Let $A = \{s_1, \dots, s_m\} \subset S^1$ be a finite augmentation
scheme with covering radius $r(A) := \sup_{s \in S^1} d(s, A)$.  Let
$\mathcal{F}$ be the class of functions $\mathcal{X} \times S^1 \to \mathbb{R}$ that
are $C$-Lipschitz in the $\tau$ argument, and suppose $f^\star \in \mathcal{F}$
is continuous-lag-equivariant and non-degenerate (i.e., its restriction
to orbits is not constant on the complement of $A$).  Then for every
learner $\hat{f}$ that uses only training data of the form
$\{(\mathrm{sh}_{s_i} h, y)\}_{s_i \in A}$,
\[
  \sup_{s \in S^1} \mathbb{E}_h \bigl|\hat{f}(h, s) - f^\star(h, s)\bigr|
  \;\ge\; C' \cdot r(A),
\]
for a constant $C' > 0$ depending on $f^\star$ but not on the learner
or the augmentation scheme.  In particular, matching an equivariant
learner's error $\varepsilon$ requires $m = \Omega(L/\varepsilon)$
augmentation points.
\end{theorem}

\begin{proof}[Proof sketch]
Pick $s^\star \in S^1$ maximizing $d(s^\star, A)$.  Choose two
histories $h_1, h_2$ whose $f^\star$-images disagree at $s^\star$ but
agree (to within Lipschitz tolerance) at every $s_i \in A$.  A learner
whose training is confined to $A$ cannot distinguish $h_1, h_2$ at
$s^\star$, so it makes an error at least $|f^\star(h_1, s^\star) -
f^\star(h_2, s^\star)|/2$.  Non-degeneracy bounds this gap below by
$C' \cdot d(s^\star, A) = C' \cdot r(A)$.  Combining with
$r(A) \ge L/(2m)$ for an $m$-point scheme gives the quantitative
$\Omega(L/m)$ lower bound.
\leanref{LEMO/ContinuousQuadrature.lean}{aug\_covering\_radius\_lower\_bound}
\end{proof}

\begin{remark}[What the theorems license in the paper's empirical
claims]\label{rem:license}
Taken together:
Theorems~\ref{thm:layer-equiv}--\ref{thm:pooled-inv} and
Lemma~\ref{lem:orbit-constancy} license the statement that LEMO is
structurally shift-equivariant/invariant at every layer and at pooled
readouts.
Theorem~\ref{thm:oodgap} and Corollary~\ref{cor:pac-gap} license the
claim that non-equivariant learners pay a $\Theta(n)$ factor in
sample complexity on shift-symmetric targets, matching our
plot of risk versus augmentation budget.
Corollaries~\ref{cor:architectural-contraction} and~\ref{cor:rollout}
license the claim that $\mathrm{LEMO}_\sigma$ has certified bounded-rollout
stability, matching the $10^{35}$-versus-$10^{-1}$ rollout gap
observed in the spectral-norm ablation.
Theorem~\ref{thm:continuous-layer-equiv}--\ref{thm:aug-lower-bound}
license the claim that only architectural lag-equivariance transfers
exactly across unseen continuous $\tau$; any finite augmentation has
error scaling at least as $\Omega(1/m)$ in the covering radius.
\end{remark}
